%% =====================================================================
%%  TEACHING DOSSIER
%% =====================================================================
%%  Options
%%    contents  -> table of contents for this document only
%%    contents depth = subsection  -> include subsections in it
%%    numbered  -> numbered sections (lettered after \appendix)
%%  Building blocks
%%    \begin{keystats}[4] \keystat{number}{label} ... \end{keystats}
%%    \ratingbar[comparison]{score}{maximum}      \ratinglegend
%%    \studentquote[source]{quotation}
%%    \begin{datatable}{column spec} ... \end{datatable}   \colhead{...}
%%        (X columns stretch; tables may break across pages)
%% =====================================================================
\begin{applicationdocument}[
  contents,
  numbered,
  subtitle = {Prepared for the position of \PositionSummary},
]{Teaching Dossier}

\tip[What a dossier is for]{A teaching dossier (called a teaching
portfolio outside Canada) pairs a reflective narrative with
\emph{evidence}. The Canadian Association of University Teachers' guide
to the teaching dossier, which most Canadian universities build on, is the
standard reference. Keep the main body to roughly 5--8 pages and move raw
material to appendices. If the posting asks only for a one- or two-page
``teaching statement'', send Section~\ref{sec:philosophy} alone.}

\section{Teaching Philosophy}\label{sec:philosophy}

I teach environmental data science because I have watched it change how
students see the world. A student who has cleaned a decade of messy lake
records and found the warming signal hidden in them never again reads a
climate headline the same way. My goal is for every student, whatever their
mathematical background, to leave my courses able to ask a question of
real environmental data, answer it honestly, and explain the
uncertainty in their answer. Three principles guide how I pursue that goal.

\subsection{Learning by doing, with real data}
Students learn data analysis by analyzing data, not by watching me do it.
In ENVS~220 I replaced two of three weekly lectures with studio sessions in
which teams work through real monitoring datasets, including records from
lakes near campus. Real data are messy. Missing values, sensor drift and
competing hypotheses are not distractions from the content; they \emph{are}
the content. On the end-of-term survey, 89\% of students agreed that the
course ``prepared me to analyze data on my own'', up from 61\% the year
before the redesign.

\subsection{Structure that includes everyone}
Inequities in prior preparation are largest in quantitative courses, and an
unstructured course quietly rewards students who already know the hidden
rules. I make expectations explicit by applying universal design for
learning and the transparent assignment framework. Every assignment states
its purpose, its task and its marking criteria. Students
receive worked examples and can choose how to present parts of their
work. A flexible ``token'' system lets anyone extend two deadlines without
explanation. These changes cost academic rigour nothing: the final project rubric
became more demanding. Yet the share of students who withdrew or failed fell
from 18\% to 9\%.

\tip[Make the philosophy specific]{For each belief, give one classroom
example and its result. Committees remember concrete scenes (``I replaced
two lectures with studios'') far better than abstractions (``I am
passionate about student-centred learning'').}

\subsection{Feedback that flows both ways}
I give students frequent, low-stakes feedback: weekly auto-graded checks
and short in-class problems discussed in pairs. I ask for the same in
return. A three-question mid-term survey (\emph{What helps? What hinders?
What would you change?}) has led to concrete changes. For example, after
students asked for them, I added recorded walkthroughs of the hardest
code. I tell the class what I changed and why, which models the
reflective practice I hope they will adopt.

\section{Teaching Responsibilities}

\begin{datatable}{@{} l X l r @{}}
  \colhead{Term} & \colhead{Course} & \colhead{Role} & \colhead{Students}\\
  \midrule\endhead
  Fall 2025   & CIVE 345 Water Quality Engineering (3rd year), Northbridge
              & Co-instructor & 54\\
  Winter 2023 & ENVS 220 Introduction to Environmental Data Analysis
                (2nd year), Eastvale & Instructor of record & 86\\
  2019--2022  & ENVS 101 Earth Systems (1st year; three terms), Eastvale
              & Teaching assistant & 168\\
  Fall 2020   & ENVS 230 Hydrology (2nd year), Eastvale & Teaching assistant & 52\\
  Winter 2021 & STAT 240 Applied Regression (2nd year), Eastvale
              & Teaching assistant & 61\\
  Fall 2021   & ENVS 410 Field Methods (4th year), Eastvale
              & Teaching assistant & 31\\
  2021--2023  & Reproducible Research with R and Git (graduate workshop
                series, 6 sessions/year) & Designer and lead & 120\\
  \bottomrule
\end{datatable}

\section{Course Design and Curriculum Development}

\subsection{Redesign of ENVS 220, Introduction to Environmental Data Analysis}
I rebuilt this required second-year course around four authentic projects
drawn from regional monitoring programs. I wrote 24 new studio
activities as accessible, captioned notebooks that run in the browser,
so students need neither a powerful laptop nor software installation. The
department has kept the redesign; two instructors have taught it since.

\subsection{Graduate workshop series}
I designed and led \emph{Reproducible Research with R and Git}, six
hands-on sessions offered each year through the School of Graduate Studies
and attended by 120 students from eleven departments. Materials are openly
licensed and have been adopted by two other institutions.

\subsection{Open educational resource}
I am writing \emph{Environmental Data Analysis in R}, a free online textbook
(CC~BY~4.0) with 14 chapters and embedded exercises. Seven chapters are
complete, and ENVS~220 and CIVE~345 already use them.

\section{Evidence of Teaching Effectiveness}

\begin{keystats}[4]
  \keystat{450+}{students taught}
  \keystat{4.6\,/\,5}{instructor rating, ENVS 220}
  \keystat{18\%\,\textrightarrow\,9\%}{withdraw-or-fail rate after redesign}
  \keystat{14}{student researchers mentored}
\end{keystats}

\tip[Triangulate your evidence]{Combine student ratings with peer
observation and with evidence of student learning (grades, products,
outcomes). Give response rates and comparison means so readers can
interpret the numbers. Quote student comments selectively but faithfully,
and never identify students.}

\subsection{Student evaluations}
The table summarizes the overall item ``The instructor was effective in
helping me learn'' (5-point scale). Item-level results are in
Appendix~\ref{app:evaluations}.

\begin{datatable}{@{} l X r l @{}}
  \colhead{Term} & \colhead{Course and role} & \colhead{Responses}
    & \colhead{My rating vs.\ comparison mean}\\
  \midrule\endhead
  F2025 & CIVE 345, co-instructor   & 39/54 (72\%) & \ratingbar[4.0]{4.5}{5}\\
  W2023 & ENVS 220, instructor      & 61/86 (71\%) & \ratingbar[4.1]{4.6}{5}\\
  F2021 & ENVS 410, TA              & 22/31 (71\%) & \ratingbar[4.3]{4.8}{5}\\
  W2021 & STAT 240, TA              & 40/61 (66\%) & \ratingbar[4.0]{4.5}{5}\\
  F2020 & ENVS 101, TA              & 47/58 (81\%) & \ratingbar[4.3]{4.7}{5}\\
  \bottomrule
\end{datatable}
\ratinglegend

\subsection{What students say}
\studentquote[ENVS 220, Winter 2023]{I came in terrified of coding and left
having built a model of a real lake. The studios made it feel possible, and
the token system meant one bad week did not sink me.}
\studentquote[CIVE 345, Fall 2025]{Dr.\ Ellis explains \emph{why} a method
works before showing how, and always connects it to a decision someone has
to make.}
\studentquote[ENVS 101, Fall 2020]{The best TA I have had. Feedback on our
lab reports was detailed and actually helped me improve.}

\subsection{Peer observation}
In March 2023, Dr.~Helen Ward of the Eastvale Centre for Teaching
Excellence observed two ENVS~220 sessions. Her report notes ``exceptionally
clear scaffolding'' and ``near-universal engagement during studio work''.
It also recommended more structured closure at the end of each studio, a
suggestion I now apply through a five-minute ``one thing I learned'' wrap-up.

\subsection{Student learning}
In the redesigned ENVS~220, the median final-project grade rose from 71\%
to 78\% under a more demanding rubric. The gap in final grades between
first-generation and continuing-generation students narrowed from 9 to 3
percentage points.

\section{Supervision and Mentorship}

I have mentored fourteen student researchers, from first-year summer
students to master's students; five are co-authors on peer-reviewed
publications. I pay undergraduate researchers rather than relying on
volunteers, so that research experience is open to students who must
work. Each mentee and I write a short mentoring plan at the start, and we
revisit it every term.

\begin{datatable}{@{} l X l @{}}
  \colhead{Years} & \colhead{Trainee and project} & \colhead{Outcome}\\
  \midrule\endhead
  2024--    & H.~Nakamura, M.Sc.: uncertainty-guided lake sampling & 1 paper, 1 talk\\
  2025--    & C.~Okonkwo, M.A.Sc.: forecasts and management decisions & 1 paper (under review)\\
  2025      & R.~Patel, NSERC USRA: LakeBench benchmark dataset & 2 papers; now Ph.D.\\
  2022--23  & W.~Chen, honours thesis: hybrid bloom models & 1 paper; now in industry\\
  2023      & \'E.~Gagnon, research assistant: volunteer monitoring data & 1 paper\\
  2020--25  & Nine further undergraduate researchers & 4 posters\\
  \bottomrule
\end{datatable}

\section{Professional Development in Teaching}
\cventry{2024}{Course Design Institute (5 days)}{Centre for Teaching and
  Learning, Northbridge Institute of Technology}{}{}
\cventry{2022}{Certificate in University Teaching (120 hours)}{Centre for
  Teaching Excellence, University of Eastvale}{}
  {Included a supervised teaching practicum and a project on inclusive
  assessment.}
\cventry{2021}{Instructor training}{The Carpentries}{}{}
\cventry{2020--2025}{Workshops on universal design for learning, inclusive
  assessment and Indigenous cultural safety}{}{}{}

\section{Teaching at \Institution: Goals and Contributions}

I am ready to teach existing courses in quantitative methods, hydrology
and environmental data analysis. I would also like to develop:
\begin{itemize}
  \item an upper-year course, \emph{Machine Learning for Environmental
        Systems}, built around projects with local partners, which would
        serve the new minor in Environmental Data Science;
  \item a graduate course, \emph{Environmental Forecasting}, covering
        process-based, statistical and hybrid forecasting with a strong
        focus on uncertainty.
\end{itemize}
Over my first five years, I aim to complete and publish the open textbook,
to measure the effect of studio-based teaching on student learning, and to
share what works through the \Faculty's teaching community.

\appendix

\section{Detailed Student Evaluation Results}\label{app:evaluations}
ENVS~220, Winter 2023 (61 of 86 students responded; scale 1--5).

\begin{datatable}{@{} X l @{}}
  \colhead{Item} & \colhead{My rating vs.\ department mean}\\
  \midrule\endhead
  The instructor explained concepts clearly.          & \ratingbar[4.0]{4.5}{5}\\
  The course was well organized.                      & \ratingbar[4.0]{4.7}{5}\\
  Assignments helped me learn.                        & \ratingbar[3.9]{4.6}{5}\\
  The instructor treated students with respect.       & \ratingbar[4.4]{4.9}{5}\\
  Feedback was timely and useful.                     & \ratingbar[3.8]{4.4}{5}\\
  Overall, the instructor was effective.              & \ratingbar[4.1]{4.6}{5}\\
  \bottomrule
\end{datatable}

\section{Sample Syllabus Excerpt: ENVS 220}
\textbf{Learning outcomes.} By the end of the course, students will be
able to:
\begin{enumerate}
  \item import, clean and document a real environmental dataset;
  \item choose and justify an appropriate statistical model;
  \item quantify and communicate uncertainty for a non-specialist audience;
  \item conduct an analysis that another person can reproduce.
\end{enumerate}

\begin{datatable}{@{} X r @{}}
  \colhead{Assessment} & \colhead{Weight}\\
  \midrule\endhead
  Weekly studio check-ins (best 10 of 12)                    & 15\%\\
  Four projects with real monitoring data                    & 50\%\\
  Peer review of a classmate's project                       & 10\%\\
  Final project: report and three-minute public explainer    & 25\%\\
  \bottomrule
\end{datatable}

\section{Supporting Documents}
The following are available on request: complete student evaluation
reports, the full peer observation report, a letter from the chair of the
Department of Environmental Science at Eastvale about the ENVS~220
redesign, and the full course materials.

\end{applicationdocument}
